\section{Proposed Architecture}
\label{sec:architecture}

\subsection{System Overview}

The target system couples a MicroBlaze~V processor, AXI interconnect, block RAM, UART, and the custom I2S transmitter IP on a Basys3 FPGA.
Audio samples are written by firmware to \texttt{DATA\_LEFT} and \texttt{DATA\_RIGHT}; the IP serializes them in Philips I\textsuperscript{2}S format to a PCM5102A module over PMOD wiring using only three signals: Bit Clock (BCLK), Word Select (WS/LRCK), and Data In (DIN).

\begin{figure}[!t]
  \centering
  \IfFileExists{processor-uart-timer-axi-i2s-architecture-whole-diagram.png}{%
    \includegraphics[width=\linewidth]{processor-uart-timer-axi-i2s-architecture-whole-diagram.png}%
  }{}
  \caption{SoC architecture with custom I2S TX IP and external PCM5102A DAC (using BCLK, WS, and DIN).}
  \label{fig:arch-soc}
\end{figure}

\subsection{Register Map and Control Fields}

Four 32-bit AXI4-Lite registers expose the peripheral to software.
\Cref{tab:regmap} lists offsets and functions; \texttt{CONTROL} packs \texttt{ENABLE}, \texttt{MUTE}, \texttt{SAMPLE\_WIDTH} (1--32 bits, 0 means 32), and \texttt{SAMPLE\_RATE\_HZ} (0--1{,}048{,}575\,Hz, safely clamped to 195.312\,kHz).

\begin{figure}[!t]
  \caption{AXI4-Lite register map of the I2S transmitter IP.}
  \label{tab:regmap}
  \centering
  \begin{tabular}{@{}llll@{}}
    \toprule
    \textbf{Offset} & \textbf{Name} & \textbf{Acc.} & \textbf{Function} \\
    \midrule
    \texttt{0x00} & \texttt{DATA\_LEFT}  & W & Left PCM sample \\
    \texttt{0x04} & \texttt{DATA\_RIGHT} & W & Right PCM sample \\
    \texttt{0x08} & \texttt{CONTROL}     & R/W & Enable, mute, width, rate \\
    \texttt{0x0C} & \texttt{STATUS}      & R & Frame-latch toggle (bit~0) \\
    \bottomrule
  \end{tabular}
\end{figure}

\subsection{DDS Clock Generation}

BCLK is derived from a phase accumulator running at $f_{\mathrm{ref}} = 100$\,MHz.
The toggle rate is $128 \times F_s$ because each BCLK half-period is generated by the accumulator wrap logic.
With a 32-bit phase path and modulus tied to 100{,}000{,}000, the average output frequency error stays below 0.001\% for standard and non-standard rates.
Although MCLK is also synthesized at $256 \times F_s$ for general compatibility, it was not required for the PCM5102A validation as the DAC was configured to generate its internal master clock from BCLK.

For stereo 32-bit slots (64 BCLK periods per frame), Philips timing requires WS transitions one BCLK before the next channel MSB and a mandatory one-bit delay after each WS edge~\cite{philips_i2s_1986}.

\subsection{Data Integrity and Handshake}

Samples written by the CPU are copied into internal shadow registers only at the start of a stereo frame (\texttt{bit\_count == 0} on a rising BCLK edge).
This prevents mid-frame corruption if software updates a data register while serialization is in progress.
\texttt{STATUS[0]} toggles whenever a new left/right pair is latched, enabling firmware to poll for the next transmit slot instead of relying on wall-clock delays or the RISC-V \texttt{mcycle} counter~\cite{riscv_isa_vol1_2019}.

\begin{figure}[!t]
  \centering
  \IfFileExists{i2s-ip-block.png}{%
    \includegraphics[width=\linewidth]{i2s-ip-block.png}%
  }{}
  \caption{Internal structure of the DDS-based I2S transmitter IP.}
  \label{fig:ip-block}
\end{figure}
